% 从旧数学附录迁出的函数图与单位边最短路说明；不新增算法。
\section{函数图的入环链与到达查询}\label{knowledge-functional-graph}
代码见第~\pageref{compact-FunctionalGraph}~页。

题目点号均从1开始，构造时统一减1；只有跳转返回的顶点需要加1，最少步数的0和不可达的-1应原样输出。
P2921只需 depth 加环长，不代表它测试了跳转或两点可达查询。
本轮完整驱动穷举至四点的全部288个后继映射，独立沿后继走到首次重复，核验步数、跳转与访问数；另验随机映射和题目最大规模。
CSES1750的题面步数上限为 $10^9$，完整unsigned64范围另由本地测试覆盖，不能由该题AC推出。

每个点有且仅有一个后继，故每个弱连通分量由一个有向环和若干指向环的树组成。
先不断移除入度为零的点，剩下的点恰在环上；沿后继枚举各环。
反向处理移除序列时，后继的环编号、入环点和入环距离均已确定，可直接传给当前点。
这里使用的是入度剥离队列，不是用显式栈模拟 DFS。

跳转 $k$ 步时，若 $k$ 小于入环距离，用普通倍增；否则先到入环点，
减去入环步数，再对环长取模。
因此即使 $k$ 达到无符号 64 位上界，也只需 $O(\log n)$ 层倍增表。
环上位置相加使用无符号宽整数，避免两个大下标的 int 中间值溢出。

判断从 $u$ 到 $v$ 的最少步数时：不同环编号必不可达。
若 $v$ 在入环树上，先检查 $u$ 的深度至少为 $v$ 的深度，
再跳过深度差，看是否恰好到 $v$，从而排除同分量的不同分支。
若 $v$ 在环上，答案为 $u$ 的入环步数加上沿后继方向的环上距离。
不可达返回 $-1$，从自己到自己为零步。
字段 position 只对环点有意义，不要将树点的 $-1$ 当作环坐标。

独立参照逐步行走直至首次重复，同时记录每个点的第一次出现位置。
穷举不超过 5 个点的全部后继映射，并检查跳转、最少步数、入环点与环长；
另验证同对象重建、空图、50 万点长链连环、50 万个自环和完整无符号 64 位步数。
普通与 ASan/UBSan 均通过。
完整驱动对应\href{https://cses.fi/problemset/task/1750}{CSES 1750}、
\href{https://cses.fi/problemset/task/1160}{CSES 1160} 和
\href{https://www.luogu.com.cn/problem/P2921}{P2921}，均使用 vector 版驱动。
前两题检查跳转与可达步数，后一题仅使用“入环步数＋环长”，不能替代查询接口的验证。
P2921 的当前程序已取得\href{https://www.luogu.com.cn/record/299973954}{线上 AC}；两道 CSES 查询题的线上评测仍待完成。

\clearpage
\section{杭州 2023 H：随机事件顺序与函数图}\label{knowledge-hangzhou-h}
\href{https://qoj.ac/problem/8237}{Sugar Sweet II} 的完整驱动为
\texttt{verify/qoj/8237.*.cpp}，复用函数图分解与逆阶乘表。
事件 $u$ 发生前，自己的数量仍为 $a_u$；目标 $v=b_u$ 只可能为 $a_v$ 或 $a_v+w_v$。
因此 $a_u<a_v$ 时必定发生，$a_u\ge a_v+w_v$ 时必定不发生。
其余情况下，事件 $u$ 成功当且仅当事件 $v$ 更早执行且成功。

将两类确定事件的后继改为自己，另保存哪些是确定成功的终点。
未修改的依赖环没有成功起点：最早执行的环上事件必失败，后续也无法形成成功链。
自环同理，自己不可能严格少于自己。
沿函数图到达确定失败终点或纯依赖环的节点，成功概率都是零。
若到达确定成功终点，且入环距离为 $d$，沿途 $d+1$ 个事件必须按唯一的顺序出现。
均匀随机全排列对这 $d+1$ 个事件诱导均匀相对顺序，因此成功概率为 $1/(d+1)!$。
这些顺序约束有关联，不能将每条依赖当作独立的二分之一概率。
期望为 $a_u+w_u/(d+1)!$，在模 $10^9+7$ 下计算。

题面所有增量为正，$\sum n\le500000$；阶乘表跨测试组只扩充不缩小，
避免大量单点测试时反复求逆。
当前驱动复用带倍增表的通用函数图，时间、空间为 $O(n\log n)$；
概率传播本身只读取 entry 和 depth，不调用跳转与最少步数接口。

验证穷举小输入及事件全排列，逐步修改真实糖果数量，最后汇总精确总和并除以排列数。
另外检查 50 万点的依赖链、纯环，以及 50 万组单点测试，普通与 ASan/UBSan 均通过。
两套驱动各 31 个正式测试与 2 个额外测试全过：
\href{https://qoj.ac/submission/2939916}{动态版} 162 ms、84608 KB，
\href{https://qoj.ac/submission/2939975}{传统版} 158 ms、84612 KB。
本应用不能验证函数图的 advance 或 steps 接口。

\newpage
\section{杭州 2023 G：开放时间的单位边最短路}\label{knowledge-hangzhou-g}
\label{app:hangzhou2023g}
\index{开放时间最短路}\index{贪吃蛇}
驱动 \texttt{verify/qoj/8236.*.cpp} 直接枚举网格邻居，不建立四倍规模的边表。
初始蛇身从头到尾编号 $0,\ldots,k-1$，第 $i\ge1$ 节所在点的最早进入时间为 $k-i$，
头与其他空格为零。旧蛇尾允许在同一步进入，所以尾格的限制必须是 1。
长度为 4 的满 $2\times2$ 网格样例距离为 $0,3,2,1$，平方和为 14。

每步只能删除或移动一个尾节，故这些时间限制必要。
最短路可以取不重复顶点的路径；行走前把途中等待的总时间全部用于缩短蛇身，
随后沿路径连续移动。进入初始蛇身第 $i$ 节时，已删除的尾节数量至少为 $k-i$；
新形成的蛇身不会与无重复顶点的路径相撞。因此这些限制也充分。
等待超过 $k-1$ 没有必要，因为此时全部初始蛇身均已释放。

转移为 $d_v=\max(d_u+1,\mathrm{open}_v)$。
按距离递增处理顶点时，第一个发现 $v$ 的前驱就给出最小候选值，故不需要后续松弛。
候选等于 $d_u+1$ 的顶点进入普通 FIFO；其余进入对应开放时间的桶。
每次取两者中距离较小的一项。桶指针只能扫描到下一个 FIFO 距离，
否则可能越过之后才被发现的较早开放时间。
令 $R=\max\mathrm{open}$，复杂度 $O(V+E+R)$，空间 $O(V+R)$。
此性质依赖所有边权为 1；一般非负权仍需优先队列。

本地以独立堆 Dijkstra 对拍，并穷举 $2\times3$ 网格障碍和全部 744 种初始蛇形，
对真实蛇身序列逐步执行移动与缩尾 BFS，核对每个点的最早到达时间。
另验证正式样例、满长链和 $3000\times3000$ 空网格及隔断网格；普通与 ASan/UBSan 均通过。
答案用 unsigned long long 累加距离平方，自然实现模 $2^{64}$。
两套驱动在 \href{https://qoj.ac/problem/8236}{QOJ8236} 各通过 58 个正式与 3 个额外测试：
\href{https://qoj.ac/submission/2940150}{动态版} 324 ms、123080 KB，
\href{https://qoj.ac/submission/2940191}{传统版} 294 ms、122768 KB。
线上证据对应单位网格、源点时间为零及初始蛇身开放时间；一般有向图与非零源点时间由本地对拍覆盖。

